% Capítulo 5 -- Citação e Referências

\section{Citação}

Citação é uma menção no texto, de informações extraídas de outras fontes \cite{ABNT2023}
e tem como objetivo dar o devido crédito ao autor de uma ideia, respeitando-se os direitos autorais.
As citações podem ser:

\begin{itemize}
  \item \textbf{Direta}: Transcrição literal de parte da obra do autor consultado. Ex.: "Deve-se indicar sempre, com método e precisão, toda documentação que serve de base para a pesquisa, assim como ideias e sugestões alheias inseridas no trabalho". \cite[p. 97]{Cervo1978}
  \item \textbf{Indireta}: Texto baseado na obra do autor consultado. Ex.: \cite{Barras1979} ressalta que, apesar da importância da arte de escrever para a ciência, inúmeros cientistas não têm recebido treinamento neste sentido.
  \item \textbf{Citação da citação}: Transcrição direta ou indireta de um texto em que não se teve acesso ao original. Ex.: "O homem é precisamente o que ainda não é. O homem não se define pelo que é, mas pelo que deseja ser" \apud{OrtegaGasset1963}[p.~160]{Salvador1977}.
\end{itemize}

As citações no texto podem seguir o sistema "numérico" ou o sistema "autor-data" \cite{ABNT2023}.

Qualquer que seja o método adotado deve ser seguido consistentemente ao longo de todo o trabalho, permitindo sua correlação na lista de referências. Não são permitidas notas de rodapé nos trabalhos do ITA.

\subsection{Sistema autor-data}

Indica-se o(s) Autor(es) pelo Sobrenome sem as iniciais, seguido do ano da publicação e da página citada, separados por vírgula.

\begin{itemize}
  \item O autor pode estar dentro ou fora dos parênteses, dependendo se for citado diretamente ou não no texto. Exemplo: Souza e Andrade (2013) estudaram um robô flexível. Robôs flexíveis apresentam graus de liberdade adicionais (Souza; Andrade, 2013). A Figura~4.2 apresenta a Estação Espacial Internacional (NASA, 2011).
  \item Quando dentro dos parênteses, os sobrenomes dos autores são separados por ponto e vírgula, por exemplo: (Autor1; Autor2; Autor3, 2007, p. 34).
  \item Citações de mais de um documento do mesmo autor. Exemplo: (Souza; Andrade, 2001, 2011, 2013).
  \item Citações de mais de um documento de autores diferentes. Exemplo: (Silva, 2003; Costa, 2004; Antunes, 2014).
  \item Quando houver coincidência de sobrenomes de autores, acrescentam-se as iniciais de seus prenomes: (Barbosa, C., 1958) e (Barbosa, O., 1958). Se mesmo assim existir coincidência, colocam-se os prenomes por extenso: (Barbosa, Cássio, 1965) e (Barbosa, Celso, 1965).
  \item As citações de diversos documentos do mesmo autor, publicados num mesmo ano, são distinguidas pelo acréscimo de letras minúsculas, em ordem alfabética, após a data e sem espacejamento. Acrescentar as letras após a data, tanto a citação, quanto na referência. Exemplo: a pesquisa apresentou um resultado (Silva, 2010a) e também outro resultado (Silva, 2010b).
\end{itemize}

\begin{table}[H]
  \centering
  \caption{Citação de Autores no texto -- ABNT NBR 10520:2023.}
  \label{tab:citacao-autores-nbr10520-2023}
  \begin{tabular}{p{0.20\textwidth}p{0.38\textwidth}p{0.38\textwidth}}
    \toprule
    \textbf{Nº de autores} & \textbf{Autor fora de parênteses} & \textbf{Autor (dentro de parênteses)} \\
    \midrule
    Um autor & Grant (1996) & (Grant, 1996) \\
    Dois autores & Zekulin e Sampaio (2010) & (Zekulin; Sampaio, 2010, p. 25) \\
    Três autores & Halliday, Resnick e Walker (2007, v. 4, p. 223) & (Halliday; Resnick; Walker, 2007, v. 4, p. 223) \\
    Mais de três autores & Takahashi et al. (1997) & (Takahashi et al., 1997, p. 8) \\
    Mesmo autor e anos diferentes & Carvalho (1986, 1993, 1996) & (Carvalho, 1986, 1993, 1996) \\
    Mesmo autor, com publicações em anos diferentes e no mesmo ano & Halkka et al. (1973, 1975a, 1975b) & (Halkka et al., 1973, 1975a, 1975b) \\
    Vários autores, mencionados simultaneamente são separados por ponto e vírgula (\;) & [... ] a mesma conclusão foi feita por Cross (2000); Knox (1986); Mezirow (2001) e Ribeiro (1989) & --- \\
    \bottomrule
  \end{tabular}
\end{table}

\subsection{Sistema numérico}

Neste sistema, a indicação da fonte é feita por uma numeração única e consecutiva, em algarismos arábicos, remetendo à lista de referências ao final do trabalho na mesma ordem em que aparecem no texto. A fonte consultada, quando repetida, deve ser representada pela mesma numeração.

No texto:

Estabilidade robusta significa que o controlador garanta estabilidade interna todos os sistemas de um grupo que inclui o sistema da planta atual$^{12}$.

Na lista de referências:

$^{12}$ OGATA, K. \textit{Engenharia de controle moderno}. Tradução I. J. Albuquerque. 2. ed. Rio de Janeiro: Prentice-Hall, 1993. 781 p.

\textit{Nota: Este template utiliza o sistema autor-data por padrão. Para utilizar o sistema numérico, é necessário alterar a opção de estilo do pacote \texttt{biblatex} em \texttt{tese.tex} de \texttt{style=abnt} para \texttt{style=abnt-numeric}.}

\section{Referências}

A Biblioteca oferece aos alunos do ITA o apoio para a normalização das referências. A solicitação desse serviço deve ser com trinta dias de antecedência da entrega da versão preliminar da tese para a Seção de Referência e Informação (BIB-RI), no e-mail bib.ref@ita.br. O arquivo com a lista de referências deverá estar no formato .docx (Word) para possibilitar correções diretamente no arquivo.

A BIB-RI retornará o arquivo com as referências padronizadas no prazo de até trinta dias úteis a partir da data do envio.

\subsection{Autoria}

Indica-se o(s) Autor(es), de modo geral, pelo último \textsc{SOBRENOME}, em maiúsculas, seguido do prenome e outros sobrenomes, abreviados ou não (apenas a letra inicial é em maiúscula).

\begin{itemize}
  \item Os autores são separados por ponto-e-vírgula, seguido de espaço.
  \item Referência com mais de 3 autores: convém indicar todos os autores na referência; permite-se que se indique apenas o primeiro autor do texto, seguido da palavra \textit{et al}.
\end{itemize}

\begin{table}[H]
\centering
\caption{Comparação de referências: manual vs.\ \texttt{\textbackslash fullcite}.}
\small
\begin{tabular}{p{0.45\textwidth}|p{0.45\textwidth}}
\toprule
\textbf{Manual} & \textbf{Biblatex (\texttt{\textbackslash fullcite})} \\
\midrule
HOWARD, N. \textit{Paradoxes of rationality}. Cambridge: MIT Press, 1971. & \fullcite{Howard1971} \\
\midrule
DEGARMO, E. P.; BLACK, J. T.; KOHSER, R. A. \textit{Materials and processes in manufacturing}. 9th ed. Hoboken: Wiley, 2003. & \fullcite{Degarmo2003} \\
\midrule
ENSSLIN, L. et al. Avaliação do desempenho de empresas terceirizadas com o uso da metodologia multicritério de apoio a decisão: construtivista. \textit{Pesquisa Operacional}, v. 30, n. 1, p. 125-152, jan./abr. 2010. & \fullcite{Ensslin2010} \\
\bottomrule
\end{tabular}
\end{table}

\begin{itemize}
  \item Organizador, editor, coordenador, compilador: quando não há autor, e sim um responsável intelectual, entra-se pelo responsável seguido da abreviação que caracteriza o tipo de responsabilidade entre parênteses e sempre no singular, independentemente do número de responsáveis: (org.), (ed.), (coord.), (comp.).
  \item Autor entidade: as obras de responsabilidade de pessoa jurídica (órgãos governamentais, empresas, associações, congressos, seminários etc.) têm entrada pela forma conhecida ou como se destaca no documento, por extenso ou abreviada (ex.: IBGE, NASA).
  \item Autoria desconhecida: entrada feita pelo título e a primeira palavra é toda em letras maiúsculas.
\end{itemize}

\fullciteblock{Waszczyszy2010}

\fullciteblock{USP1993}

\fullciteblock{ITA1958}

\fullciteblock{DiagnosticoEditorial1993}

\subsection{Título e subtítulo}

O título e o subtítulo (se for usado) devem ser reproduzidos tal como figuram no documento, separados por dois-pontos. O subtítulo nas referências não leva destaque.

Em títulos e subtítulos demasiadamente longos, podem-se suprimir as últimas palavras, desde que não seja alterado o sentido. A supressão deve ser indicada por reticências entre colchetes, por exemplo: [\ldots]. Quando o título aparecer em mais de uma língua, registra-se o primeiro. Opcionalmente, registra-se o segundo ou o que estiver em destaque, separando o primeiro pelo sinal de igualdade. Os títulos dos periódicos podem ser abreviados, desde que conste na publicação.

Exemplos:

\fullciteblock{Bitter2001}

\fullciteblock{Tedlow2012}

\fullciteblock{SaoPauloMedicalJournal}

\fullciteblock{Gray2012}

\subsection{Edição}

Quando houver uma indicação de edição, esta deve ser transcrita, utilizando-se abreviaturas dos numerais ordinais e da palavra edição, ambas de forma adotada no idioma do documento. Ex.: 2. ed. (português e espanhol), 2nd ed. (inglês), 2. Aufl. (alemão), 2e ed. (francês) e 2a ed. (italiano). Considerar a versão de documentos eletrônicos como equivalente à edição e transcrevê-la como consta no documento.

\fullciteblock{Astrology1994}

\subsection{Imprenta (Local: editora, data)}

O nome do local (cidade) de publicação deve ser indicado tal como figura no documento. Na ausência do nome da cidade, pode ser indicado o estado ou o país, desde que conste no documento. No caso de homônimos de cidades, acrescenta-se a sigla do estado, do país etc. Ex.: Viçosa, AL / Viçosa, MG / Viçosa, RJ. Quando a cidade não aparece no documento, mas pode ser identificada, indica-se entre colchetes.

\fullciteblock{Ogata1993}

\fullciteblock{LazzariniNeto1994}

Não sendo possível determinar o local, utiliza-se a expressão \textit{sine loco}, abreviada entre colchetes e em itálico: [S.\ l.].

\subsection{Editora}

O nome da Editora ou outra instituição responsável pela publicação deve ser indicado tal como figura no documento, suprimindo-se as palavras que designam a natureza jurídica ou comercial. Quando a editora não puder ser identificada, deve-se indicar expressão \textit{sine nomine}, abreviada, entre colchetes e em itálico: [s.\ n.]. Quando o local e a editora não puderem ser identificados na publicação, utilizam-se ambas as expressões, abreviadas e entre colchetes: [S.\ l.: s.\ n.].

\subsection{Data}

A data é um elemento obrigatório para a referência; sempre deve ser identificada seja da publicação, distribuição, do copirraite, da impressão, da apresentação (depósito) de um trabalho acadêmico, ou outra.

Não identificando nenhuma data de publicação, distribuição, copirraite, impressão etc., deve ser indicado um ano entre colchetes, conforme indicado:

\begin{itemize}
  \item \mbox{[1971 ou 1972]} um ano ou outro;
  \item \mbox{[1969?]} ano provável;
  \item \mbox{[1973]} ano certo, não indicado no item;
  \item \mbox{[entre 1906 e 1912]} usar intervalos menores de 20 anos;
  \item \mbox{[ca. 1960]} ano aproximado;
  \item \mbox{[197-]} década certa;
  \item \mbox{[197-?]} década provável;
  \item \mbox{[18--]} século certo;
  \item \mbox{[18--?]} século provável.
\end{itemize}

\subsection{Descrição física}

Quanto à paginação, opcionalmente registrar o número total de páginas (ex.: 530 p.). Para referenciar partes de publicações, deve-se indicar os números das folhas ou páginas inicial e final, precedidos da abreviatura f. ou p.; ou indicar o número do volume, precedido de v.; ou indicar o número do capítulo, precedido de cap. Se necessário, indicar outra forma de individualizar a parte referenciada.

Ex.: p. 56--59; cap. 5; v. 1, p. 20--25. Quando o documento for publicado em mais de uma unidade física, ou seja, mais de um volume, indica-se a quantidade de volumes, seguida da abreviatura v.

\fullciteblock{Nussenzveig2002}

\subsection{Trabalhos acadêmicos}

Nas teses, dissertações ou outros trabalhos acadêmicos deve ser indicado o tipo de documento (tese, dissertação, trabalho de conclusão de curso, mestrado profissional etc.), o grau, a vinculação acadêmica, o local e a data da defesa, mencionada na folha de aprovação (se houver).

\fullciteblock{CardosoJunior2014}

\section{Ordenação das referências}

As referências devem ser elaboradas em espaço simples, alinhadas à margem esquerda do texto e separadas entre si por uma linha em branco de espaço simples. As referências dos documentos citados em um trabalho devem ser ordenadas no sistema numérico quando o sistema de citação for numérico ou por ordem alfabética quando o sistema for autor-data. Segue abaixo alguns exemplos de referências bibliográficas de acordo com o tipo de fonte.

\subsection{Artigo de periódico}

AUTOR. Título do artigo. Título do Periódico/``Source'', v. (volume), n. (número) do fascículo, página inicial--página final, mês ano.

\fullciteblock{SusanResiga1988}

\fullciteblock{Cavalcante2000}

\fullciteblock{Denton1999}

\subsection{Eventos: trabalhos apresentados em congressos, seminários, simpósios, workshop}

AUTOR do ARTIGO. Título do artigo: subtítulo. In: NOME DO EVENTO, número, ano, local de realização. Anais [\ldots] (\textit{Proceedings} [\ldots] quando em inglês). Local de publicação (cidade): Editora, data. volume, páginas inicial-final do trabalho, se houver.

\fullciteblock{Bai2001}

BENITTI, F. B. V.; ZIMMERMANN, A. P. Testware: ferramenta de planejamento e execução de casos de teste. In: SEMINÁRIO DE COMPUTAÇÃO, 2007, Blumenau. Anais [\ldots]. Blumenau: Universidade Regional de Blumenau, 2007. p. 63-74.

\fullciteblock{Benitti2007}

\fullciteblock{Luchinsky2009}

\fullciteblock{Gimenes2008}

\fullciteblock{Candido1997}

\fullciteblock{Ribeiro2010}

\subsection{Livros e folhetos}

Os elementos essenciais de uma monografia no todo são: autor(es), título, subtítulo (se houver), edição (se houver), local, editora e data de publicação.

AUTOR. Título: subtítulo. Edição. Local (cidade) de publicação: Editora, data. Número de páginas ou volumes. (Título da Série, volume da série).

\fullciteblock{Franco2004}

\fullciteblock{Knuth1999}

\subsection{Capítulo de livro}

AUTOR DO CAPÍTULO. Título do capítulo. In: AUTOR/EDITOR DO LIVRO. Título: subtítulo do livro. Edição (se houver). Local (Cidade) de publicação: Editora, data. volume, capítulo, páginas inicial-final da parte.

\fullciteblock{Zilio2001}

\fullciteblock{Perrins1983}

\subsection{Documentos de acesso exclusivo em meio eletrônico}

AUTORIA. Título do serviço ou produto. Versão (se houver). Local (cidade) de publicação: editor, data de publicação. Disponível em: indicar endereço eletrônico. Acesso em: dia, mês e ano.

Nota: no caso de arquivos eletrônicos, acrescentar a respectiva extensão à denominação atribuída ao arquivo.

Obs.: para qualquer tipo de referência, as informações acrescidas devem seguir o idioma do texto do trabalho e não do documento referenciado; ou seja, utilize \textit{Available at:} e \textit{Accessed on:} apenas no caso de estar redigindo o seu trabalho acadêmico em inglês.

\fullciteblock{Gaines1995}

\fullciteblock{ForecastInternational2006}

\subsection{Trabalhos acadêmicos, monografias, teses, dissertações, TG}

AUTOR. Título: subtítulo (se houver). Ano do depósito. Tipo de trabalho (Grau e Curso) -- Nome da Universidade, Cidade, data de apresentação ou defesa. Número de páginas (opcional).

\fullciteblock{Naves2006}

\fullciteblock{Mizioka2009}

\fullciteblock{Sharma2006}

\subsection{Artigo de jornal}

AUTOR. Título do artigo e subtítulo (se houver). Título do Jornal, subtítulo do jornal (se houver). Local da publicação, numeração do ano e/ou volume, número (se houver), data de publicação, seção, caderno ou parte do jornal e a paginação correspondente. Quando não houver seção, caderno ou parte, a paginação do artigo ou matéria precede a data.

\fullciteblock{Leal1999}

\fullciteblock{Azevedo1985}

\subsection{Normas técnicas}

AUTORIDADE. Número da norma: título e subtítulo. Local da publicação (cidade): Editor, data. Número de páginas (opcional).

\fullciteblock{MIL1797}

\fullciteblock{ABNTNBR14724}

\fullciteblock{ICAO2007}

\subsection{Patentes}

INVENTOR (Autor). Título da invenção na língua original. Nome do depositante e do procurador (se houver). Número da patente, datas (de depósito e de concessão, se houver). Indicação da publicação onde foi citada a patente, quando for o caso.

\fullciteblock{Wischoff1995}

\fullciteblock{Bertazzoli2006}

\subsection{Programa de computador}

AUTOR. Título: subtítulo (se houver). Versão. Local de publicação: Editora, data da publicação. Tipo de mídia. Descrição física. Notas opcionalmente para melhor explicar o material consultado.

\fullciteblock{Microsoft1995}

\subsection{Relatório técnico}

AUTOR. Título: subtítulo (se houver). Edição. Local (cidade) de publicação: Editora, data. (Série, v.).

\fullciteblock{Allison1993}

\fullciteblock{Dunham1998}

\subsection{Título do periódico}

TÍTULO DO PERIÓDICO. Local de publicação (cidade): Editora, datas de início e de encerramento da publicação (se houver), e ISSN (se houver).

\fullciteblock{RevistaBrasileiraGeografia}

\fullciteblock{SaoPauloMedicalJournal}

\subsection{Citação de citação}

Todo esforço deve ser empreendido para se consultar o documento original. Entretanto, nem sempre é possível o acesso a certos textos. Nesse caso, apesar de não ser a maneira recomendada, pode-se reproduzir informação já citada por outros autores, cujos documentos tenham sido efetivamente consultados. Pode-se adotar o seguinte procedimento:

No texto, citar o sobrenome do autor do documento não consultado, seguido das expressões: citado por, \textit{apud}, conforme ou segundo, e o sobrenome do autor do documento efetivamente consultado. Na listagem de referência, elencar somente a fonte consultada.

Exemplo:

No texto: Marinho (citado por Marconi e Lakatos, 1982) apresenta a formulação do problema como uma fase de pesquisa que, sendo bem delimitado, simplifica e facilita a maneira de conduzir a investigação.

Na lista de referências:

\fullciteblock{MarconiLakatos1982}
